\subsection{Calculation of Gaussian integrals}

Lemma 1. --- For every integer \(n \geqslant 0\) ,

\begin{enumerate}
\def\labelenumi{(\arabic{enumi})}
\setcounter{enumi}{5}
\tightlist
\item
\end{enumerate}

\[
\int_ {\mathbf {R}} | x | ^ {n} e ^ {- x ^ {2} / 2} d x = 2 ^ {\frac {n + 1}{2}} \Gamma \left(\frac {n + 1}{2}\right)\tag{7}
\]

\[
\int_ {\mathbf {R}} x ^ {2 n} e ^ {- x ^ {2} / 2} d x = (2 \pi) ^ {1 / 2} \frac {(2 n) !}{2 ^ {n} n !}\tag{8}
\]

\[
\int_ {\mathbf {R}} x ^ {2 n + 1} e ^ {- x ^ {2} / 2} d x = 0.
\]

Recall the formula

\[
\Gamma (s) = \int_ {0} ^ {\infty} u ^ {s - 1} e ^ {- u} d u\tag{9}
\]

valid for every real number \(s > 0\) (FRV, VII, §1, No.~3, Prop. 3). On making the change of variable \(x = (2u)^{1/2}\) , it follows from (9) that

\[
\int_ {0} ^ {\infty} x ^ {n} e ^ {- x ^ {2} / 2} d x = \int_ {0} ^ {\infty} (2 u) ^ {n / 2} e ^ {- u} \frac {1}{2} 2 ^ {1 / 2} u ^ {- 1 / 2} d u = 2 ^ {\frac {n - 1}{2}} \Gamma \left(\frac {n + 1}{2}\right),
\]

whence the formula (6) since

\[
\int_ {\mathbf {R}} | x | ^ {n} e ^ {- x ^ {2} / 2} d x = 2 \int_ {0} ^ {\infty} x ^ {n} e ^ {- x ^ {2} / 2} d x.
\]

The formula (7) follows from (6) and the relation

\[
\Gamma \left(n + \frac {1}{2}\right) = \pi^ {1 / 2} \frac {(2 n) !}{2 ^ {2 n} n !}.\tag{10}
\]

For \(n = 0\) , this relation reduces to \(\Gamma(\frac{1}{2}) = \pi^{1/2}\) , that is, to the formula (21) of FRV, VII, §1, No.~3. The general case then follows by induction on \(n\) , on taking into account the relation \(\Gamma(x + 1) = x \cdot \Gamma(x)\) (loc. cit., §1, No.~1).

Finally, the formula (8) follows from the fact that the function \(x \mapsto x^{2n+1}e^{-x^2/2}\) is odd.

Lemma 2. --- For every complex number y,

\[
(2 \pi) ^ {- 1 / 2} \int_ {\mathbf {R}} e ^ {- x ^ {2} / 2} e ^ {i x y} d x = e ^ {- y ^ {2} / 2}.\tag{11}
\]

In particular,

\[
(2 \pi) ^ {- 1 / 2} \int_ {\mathbf {R}} e ^ {- x ^ {2} / 2} d x = 1.
\]

The change of variable \(x \mapsto -x\) yields

\[
(2 \pi) ^ {- 1 / 2} \int_ {\mathbf {R}} e ^ {- x ^ {2} / 2} e ^ {i x y} d x = (2 \pi) ^ {- 1 / 2} \int_ {\mathbf {R}} e ^ {- x ^ {2} / 2} e ^ {- i x y} d x;
\]

since \(\cos u = \frac{e^{iu} + e^{-iu}}{2}\) for every complex number \(u\) , it follows that

\[
(2 \pi) ^ {- 1 / 2} \int_ {\mathbf {R}} e ^ {- x ^ {2} / 2} e ^ {i x y} d x = (2 \pi) ^ {- 1 / 2} \int_ {\mathbf {R}} e ^ {- x ^ {2} / 2} \cos x y d x.\tag{12}
\]

For every integer \(n \geqslant 0\) , set

\[
g _ {n} (x) = (- 1) ^ {n} (2 \pi) ^ {- 1 / 2} \frac {(x y) ^ {2 n}}{(2 n) !} e ^ {- x ^ {2} / 2}.
\]

By (7),

\begin{enumerate}
\def\labelenumi{(\arabic{enumi})}
\setcounter{enumi}{12}
\tightlist
\item
\end{enumerate}

\[
\int_ {\mathbf {R}} | g _ {n} (x) | d x = \frac {1}{n !} \biggl (\frac {| y | ^ {2}}{2} \biggr) ^ {n}\tag{14}
\]

\[
\int_ {\mathbf {R}} g _ {n} (x) d x = \frac {1}{n !} \Big (- \frac {y ^ {2}}{2} \Big) ^ {n},
\]

whence

\[
\sum_ {n = 0} ^ {\infty} \int_ {\mathbf {R}} | g _ {n} (x) | d x = e ^ {| y | ^ {2} / 2} <   + \infty .
\]

Since, moreover,

\[
(2 \pi) ^ {- 1 / 2} e ^ {- x ^ {2} / 2} \cos x y = \sum_ {n = 0} ^ {\infty} g _ {n} (x),
\]

this equality can be integrated term by term, whence

\[
(2 \pi) ^ {- 1 / 2} \int_ {\mathbf {R}} e ^ {- x ^ {2} / 2} \cos x y d x = \sum_ {n = 0} ^ {\infty} \int_ {\mathbf {R}} g _ {n} (x) d x = e ^ {- y ^ {2} / 2}
\]

by (14). The formula (11) then follows from (12).

